\naslov{Sistemi linearnih enačb}

\begin{definicija}
  Preslikava $\norm{\cdot} : \C^n \to \R$ je \pojem{vektorska norma}, če velja
  \begin{itemize}
  \item $\norm{x} \ge 0$ za vse $x$, in $\norm{x} = 0 \nt x = 0$,
  \item $\norm{\alpha x} = \abs{\alpha} \norm{x}$,
  \item $\norm{x + y} \le \norm{x} + \norm{y}$.
  \end{itemize}
\end{definicija}

Vse vektorske norme so ekvivalentne, za poljubni normi $\norm{x}_A$ in
$\norm{x}_B$ obstajata konstanti $C_1, C_2$, da velja
\[
  C_1 \norm{x}_A \le \norm{x}_B \le C_2 \norm{x}_A.
\]
Konkretno za 1-normo, 2-normo in supremum normo veljajo ocene
\begin{gather*}
  \norm{x}_2 \le \norm{x}_1 \le \sqrt{n} \norm{x}_2 \\
  \norm{x}_\infty \le \norm{x}_1 \le \sqrt{n} \norm{x}_\infty \\
  \norm{x}_\infty \le \norm{x}_2 \le \sqrt{n} \norm{x}_\infty
\end{gather*}

\vprasanje{Definiraj vektorsko normo. V kakšnem razmerju so znane vektorske
  norme?}

\begin{definicija}
  Preslikava $\norm{\cdot} : \C^{n \times n} \to \R$ je \pojem{matrična norma},
  če velja
  \begin{itemize}
  \item $\norm{A} \ge 0$ in $\norm{A} = 0 \nt A = 0$,
  \item $\norm{\alpha A} = \abs{\alpha} \norm{A}$,
  \item $\norm{A + B} \le \norm{A} + \norm{B}$,
  \item $\norm{A B} \le \norm{A} \norm{B}$ (submultiplikativnost).
  \end{itemize}
\end{definicija}

Matrika je tudi vektor, na njej so tudi definirane običajne vektorske norme.
Definiramo funkcije
\begin{gather*}
  N_1(A) = \sum_{i,j} \abs{a_{ij}}, \\
  N_2(A) = \sqrt{\sum_{i,j} \abs{a_{ij}}^2}, \\
  N_\infty(A) = \max_{i,j} \abs{a_{ij}}.
\end{gather*}
Te funkcije ustrezajo prvim trem točkam definicije matrične norme, to pa še ni
dovolj.
Za matriki
\[
  A = B =
  \begin{bmatrix}
	1 & 1 \\ 1 & 1
  \end{bmatrix}
\]
velja $N_\infty(A) = N_\infty(B) = 1$, ampak $N_\infty(AB) = 2$, torej
$N_\infty$ ni matrična norma.
V nasprotju $N_1$ in $N_2$ dejansko sta matrični normi.
Funkcijo $N_2$ imenujemo \pojem{Frobeniusova norma} in označimo z $N_2(A) =
\norm{A}_F$.

\vprasanje{Dokaži, da $N_\infty$ ni matrična norma. Kaj je Frobeniusova norma?}

\begin{izrek}
  Naj bo $\norm{\cdot}_v$ vektorska norma na $\C^n$.
  Potem je
  \[
	\norm{A}_m = \max_{x \ne 0} \frac{\norm{A x}_v}{\norm{x}_v}
  \]
  matrična norma.
  Taki normi pravimo \pojem{operatorska norma}.
\end{izrek}

\begin{proof}
  Prve tri točke so očitne, preverimo samo submultiplikativnost.
  Za vsak $x \ne 0$ velja $\norm{A x}_v \le \norm{A}_m \norm{x}_v$ po
  definiciji, torej
  \[
	\norm{AB}
	= \max_{x \ne 0} \frac{\norm{ABx}_v}{\norm{x}_v}
	\le \max_{x \ne 0} \frac{\norm{A}_m \norm{Bx}_v}{\norm{x}_v}
	= \norm{A}_m \norm{B}_m.
  \]
\end{proof}

\vprasanje{Dokaži, da so operatorske norme res matrične norme.}

Za matrično normo $\norm{\cdot}_m$ in vektorsko normo $\norm{\cdot}_v$ pravimo,
da sta \pojem{usklajeni}, če za vsako matriko $A$ in vektor $x$ velja
\[
  \norm{Ax}_v \le \norm{A}_m \norm{x}_v.
\]

\begin{lema}
  Za vsako matrično normo obstaja usklajena vektorska norma.
\end{lema}

\begin{proof}
  Za vektor $x$ definiramo
  \[
	\norm{x}_v =
	\norm{
	  \begin{bmatrix}
		x & 0 & \ldots & 0
	  \end{bmatrix}
	}_m,
  \]
  kjer smo vektor dopolnili do kvadratne matrike.
  To je očitno vektorska norma, normi sta očitno usklajeni.
\end{proof}

\vprasanje{Dokaži, da ima vsaka matrična norma usklajeno vektorsko normo.}

\begin{posledica}
  Za vsako matrično normo in poljubno lastno vrednost $\lambda$ matrike $A$
  velja $\abs{\lambda} \le \norm{A}$.
\end{posledica}

\begin{proof}
  Naj bo $Ax = \lambda x$ za nek $x \ne 0$.
  Velja
  \[
	\abs{\lambda} \norm{x}_v = \norm{\lambda x}_v = \norm{Ax}_v \le \norm{A}
	\norm{x}_v.
  \]
\end{proof}

\vprasanje{Kakšna je povezava med lastnimi vrednostmi in matrično normo?
  Dokaži.}

\begin{lema}
  Norma $\norm{A}_1$ je enaka največji 1-normi stolpca matrike $A$.
\end{lema}

\begin{proof}
  Naj bo $x$ vektor.
  Velja
  \[
	\norm{Ax}_1
	= \norm{\sum_i x_i a_i}_1
	\le \sum_i \abs{x_i} \norm{a_i}_1
	\le \max_{j=1, \ldots, n} \sum_i \abs{x_i} \norm{a_j}_1
	= \max_j \norm{x}_1 \norm{a_j}_1.
  \]
  Sledi
  \[
	\frac{\norm{Ax}_1}{\norm{x}_1} \le \max_j \norm{a_j}_1.
  \]
  Enakost dobimo, če za $x$ vzamemo $e_k$, kjer je $k$ stolpec z največjo
  1-normo.
\end{proof}

Podobno pokažemo, da je $\norm{A}_\infty$ enaka največji 1-normi vrstice.

\vprasanje{Čemu sta enaki normi $\norm{A}_1$ in $\norm{A}_\infty$? Dokaži za
  eno.}

\begin{lema}
  Velja $\norm{A}_2 = \max_j \sqrt{\lambda_j}$, kjer so $\lambda_j$ lastne
  vrednosti matrike $A^H A$.
\end{lema}

\begin{proof}
  Matrika $A^H A$ je hermitska, torej so vse njene lastne vrednosti realne.
  Velja
  \[
	x^H A^H A x = (Ax)^H Ax = \norm{Ax}_2 \ge 0,
  \]
  torej je $A^H A$ pozitivno semidefinitna, in ima nenegativne lastne vrednosti.
  Izraz je torej res dobro definiran.

  Naj bodo $\sigma_1^2 \le \sigma_2^2 \le \ldots \le \sigma_n^2$ singularne
  vrednosti matrike $A$ (lastne vrednosti $A^H A$).
  Ker je $A^H A$ hermitska, lahko lastne vektorje izberemo tako, da so
  ortonormirani.
  Naj za $v_i$ torej velja $A^H A v_i = \sigma_i^2 v_i$.
  Za poljuben $x$ velja
  \[
	x = \sum_i \alpha_i v_i,
  \]
  torej
  \[
	\norm{Ax}_2^2
	= (Ax)^H (Ax)
	= x^H A^H A x
	= \sum_i \abs{\alpha_i}^2 \sigma_i^2
	\le \sigma_n^2 \sum_i \abs{\alpha_i}^2
	= \sigma_n^2 \norm{x}_2^2.
  \]
  Sledi neenakost
  \[
	\frac{\norm{Ax}_2}{\norm{x}_2} \le \sigma_n,
  \]
  kjer dobimo enačaj, če vzamemo $x = v_n$.
\end{proof}

\vprasanje{Čemu je enaka 2-norma matrike? Dokaži.}

